\chapter{从有限定理走向更强理论}
\label{ch:frontiers}
\chapterguide{进一步研究应当改进核心模型的可达范围与解释力，并保留明确的反例和比较边界。}

\section{当前已经连接起来的结果}
不可辨别类决定永久误差，覆盖率决定一个可执行组织的可靠性，自适应KL账本约束全部策略。在固定$d,M$下，三者形成操作性主定理。校准模型给出表示存在性与构造算法；通信后处理、充分计数和量化界说明组织改变何时保存这些保证。

这条链采用经典统计实验、主动检验、信息不等式和有限线性代数。其价值应根据模型是否清楚、连接是否正确、解释是否可复用来判断。哪些结果能形成独立原创贡献，还需要逐项与最接近的研究比较；此处没有把重新命名的经典量列作发现。

\section{问题一：规模增长时，覆盖率还缺少什么}
主定理含$d$因子，上界含$\ln m$。当情境数随问题规模指数增长时，这些因素可能压倒结论。需要研究目标冲突图、概率表的对称性和可构造表示，能否替代原始情境数。

二元模型提供一个具体测试：标准基实验与$k$位异或目标具有$\mu=1/k$，局部多数构造需要$O(k\ln(k/\delta)/(1/2-\eta)^2)$次调用，而通用最大似然并合界可能更松。能否对完整自适应类别证明更紧下界，并找出决定额外对数因子的结构？这里可以检验覆盖率是否需要新的互补参数。

\section{问题二：什么时候适应性具有不可替代收益}
全局覆盖分配同时保护全部情境，局部KL规划允许对真实情境分别设计。要研究两者之间的实际差距，需要把学习当前所在区域的成本加入策略，并比较最坏风险、平均风险、固定预算和期望停止时间等不同标准。

一个有内容的分离应同时给出：所有非自适应组织的下界、可执行自适应组织的上界，以及不依赖新增观察权限的比较。不能仅凭两个形式不同的规划值就宣布可达的适应性优势。

\section{问题三：共同错误在过程中变化}
静态$\theta$表达持续的共同因素。若一次交谈改变其他成员的错误状态，或一次操作破坏未来观察条件，输出核需要依赖隐藏状态与历史。组织此时同时选择取证与改变情境。

可以先限定有限状态、已知转移和局部可见性，研究是否存在任何能把两个目标冲突状态区分开的控制序列。静态模型中的一次实验混淆类将变成受控可辨别性问题。新的上下界应把等待、干预和错误传播计入状态转移。

\section{问题四：观察模型怎样取得}
已知概率表把研究集中在组织。未知表需要从任务与诊断样本中学习，并保证对小概率项、共享偏差和未见情境的处理。第\ref{thm:robust}条定理只说明统一对数误差如何消耗余量，尚未提供取得该误差保证的通用算法。

一个可推进的任务是：固定具有有限描述的概率表族，联合控制模型识别成本与之后的诊断成本，寻找先学习再组织与边学边组织的分界。允许任意未知表时，未观察区域上的不同解释可能永远不可辨别。

\section{问题五：通信怎样内生改变组织}
完整计数可以保存自适应组织的全部证据，选择性通信则未必交付完整计数。私有信息驱动的消息既改变接收者后验，又影响其他成员下一步实验。联合优化需要同时描述消息生成、可见沉默与局部策略，并对全部相关成本给出界。

可先研究两个成员、有限观察与有代价的可靠时隙：比较固定轮询、接收者请求和事件触发的最优任务风险。若某种协议取得优势，需要证明优势来自更有效地利用已有局部信息，而未在发送时刻或控制信号中隐藏免费信道。

\section{研究成果的检验单位}
一项结果应能独立交付模型、定理、证明、反例与最接近文献的差异。对一般性结论，数值核验只负责发现有限实例中的错误；对实际Agent，则需另外估计实验模型并检验失配。

可以优先解决上述某一个窄问题，再看新结构是否能回过头改进主定理、校准算法和通信分析。由同一结果推动多个部分，才能使理论继续统一，而非依靠增加概念扩大篇幅。

\begin{exercise}选择一个增长的有限实例族，指出主定理中哪个常数随规模变化，并说明它是否仍给出有用界。\end{exercise}
\begin{exercise}为一个候选研究结果列出最小可审查材料：模型、上下界、反例、可复算实例和最近已有结果。\end{exercise}
